\documentclass{article}
\usepackage{amsmath,amssymb,amsthm}
\usepackage{algorithm,algorithmic}
\usepackage{hyperref}
\newtheorem{assumption}{Assumption}
\newtheorem{theorem}{Theorem}
\newtheorem{lemma}{Lemma}
\newtheorem{corollary}{Corollary}
\begin{document}

\section*{Cyclic Stochastic Gradient Descent (Cyclic‑SGD)}

\section*{Mathematical Formulation}
Let $f:\mathbb{R}^{d}\rightarrow\mathbb{R}$ be a (possibly non‑convex) loss function.  
At iteration $t\in\{0,1,\dots,T-1\}$ we draw an i.i.d. random sample $z_t$ and compute a stochastic gradient
\[
g_t \;:=\;\nabla f(w_t;z_t)\in\mathbb{R}^{d}.
\]
The learning rate follows a *triangular cyclic schedule* of period $M\in\mathbb{N}_{+}$:
\[
\eta_t \;=\;
\begin{cases}
\displaystyle \eta_{\min}+ \frac{t\bmod M}{M/2}\,(\eta_{\max}-\eta_{\min}), & 0\le (t\bmod M)< M/2,\\[2ex]
\displaystyle \eta_{\max}- \frac{t\bmod M-M/2}{M/2}\,(\eta_{\max}-\eta_{\min}), & M/2\le (t\bmod M)< M,
\end{cases}
\tag{1}
\]
where $0<\eta_{\min}<\eta_{\max}$ are the minimum and maximum step sizes.  
The update rule is
\[
\boxed{\,w_{t+1}= w_t-\eta_t\, g_t\, } \tag{2}
\]

\section*{Pseudocode}
\begin{algorithm}[H]
\caption{Cyclic‑SGD}
\begin{algorithmic}[1]
\Require Initial point $w_{0}\in\mathbb{R}^{d}$, cycle length $M$, bounds $0<\eta_{\min}<\eta_{\max}$, total iterations $T$.
\For{$t=0$ \textbf{to} $T-1$}
    \State Sample $z_t\sim\mathcal{D}$ and compute $g_t=\nabla f(w_t;z_t)$
    \State Compute $\eta_t$ by (1)   \Comment{triangular schedule}
    \State $w_{t+1}= w_t-\eta_t g_t$   \Comment{SGD step}
\EndFor
\State \textbf{return} $w_T$
\end{algorithmic}
\end{algorithm}

\section*{Dry Run Example}
Consider the univariate quadratic $f(w)=(w-3)^2$, whose exact gradient is $\nabla f(w)=2(w-3)$.  
We take a deterministic ``stochastic’’ gradient $g_t=\nabla f(w_t)$, $w_0=0$, $M=4$, $\eta_{\min}=0.05$, $\eta_{\max}=0.15$ and run three iterations.

\begin{center}
\begin{tabular}{c|c|c|c}
$t$ & $\eta_t$ (by (1)) & $g_t=2(w_t-3)$ & $w_{t+1}=w_t-\eta_t g_t$ \\ \hline
0 & $0.05$ & $2(0-3)=-6$ & $0-0.05(-6)=0.30$ \\
1 & $0.10$ & $2(0.30-3)=-5.40$ & $0.30-0.10(-5.40)=0.84$ \\
2 & $0.15$ & $2(0.84-3)=-4.32$ & $0.84-0.15(-4.32)=1.496$ 
\end{tabular}
\end{center}
Corresponding objective values: $f(0)=9$, $f(0.30)=7.29$, $f(0.84)=4.62$, $f(1.496)=2.26$.

\newpage
\section*{Assumptions}
\begin{assumption}[Smoothness]\label{ass:smooth}
$f$ is $L$‑smooth, i.e. for all $w,u\in\mathbb{R}^{d}$
\[
\|\nabla f(w)-\nabla f(u)\|\le L\|w-u\| .
\tag{3}
\]
Equivalently,
\[
f(u)\le f(w)+\langle\nabla f(w),u-w\rangle+\frac{L}{2}\|u-w\|^{2}.
\tag{4}
\]
\end{assumption}

\begin{assumption}[Unbiased Stochastic Gradients]\label{ass:unbiased}
For every $t$, conditional on $w_t$,
\[
\mathbb{E}[g_t\mid w_t]=\nabla f(w_t).
\tag{5}
\]
\end{assumption}

\begin{assumption}[Bounded Variance]\label{ass:variance}
There exists $\sigma^{2}\ge0$ such that
\[
\mathbb{E}\big[\|g_t-\nabla f(w_t)\|^{2}\mid w_t\big]\le\sigma^{2}\qquad\forall t .
\tag{6}
\]
\end{assumption}

\begin{assumption}[Lower Bounded Objective]\label{ass:lower}
$f$ is bounded below: $f^{*}:=\inf_{w}f(w)>-\infty$.
\tag{7}
\end{assumption}

\section*{Main Convergence Theorem}
\begin{theorem}[Convergence of Cyclic‑SGD]\label{thm:main}
Assume \ref{ass:smooth}–\ref{ass:lower} hold and let $\{\eta_t\}_{t=0}^{T-1}$ be generated by the triangular schedule (1) with
\[
0<\eta_{\min}\le\eta_t\le\eta_{\max}\qquad\forall t .
\tag{8}
\]
Define the averaged iterate $\bar w_T:=\frac{1}{T}\sum_{t=0}^{T-1}w_t$.  
Then
\[
\frac{1}{T}\sum_{t=0}^{T-1}\mathbb{E}\!\big[\|\nabla f(w_t)\|^{2}\big]
\;\le\;
\underbrace{\frac{2\big(f(w_{0})-f^{*}\big)}{\eta_{\min}T}}_{\text{bias term}}
\;+\;
\underbrace{L\,\eta_{\max}\,\sigma^{2}}_{\text{variance neighbourhood}} .
\tag{9}
\]
Consequences:
\begin{enumerate}
\item If $\eta_{\max}=O\!\big(\tfrac{1}{\sqrt{T}}\big)$ (e.g. decrease both $\eta_{\min}$ and $\eta_{\max}$ at that rate) then
\[
\frac{1}{T}\sum_{t=0}^{T-1}\mathbb{E}\!\big[\|\nabla f(w_t)\|^{2}\big]=O\!\big(T^{-1/2}\big).
\]
\item With a *fixed* cyclic schedule (constant $\eta_{\min},\eta_{\max}$) the iterates converge to a neighbourhood of a stationary point whose radius is $O(L\eta_{\max}\sigma^{2})$.
\end{enumerate}
\end{theorem}

\section*{Theoretical Properties}
\begin{itemize}
\item \textbf{Stability.} The upper bound (9) shows that the algorithm is stable as long as the maximal step size $\eta_{\max}$ is chosen such that $L\eta_{\max}<1$; otherwise the $L\eta_{\max}\sigma^{2}$ term may dominate and the bound becomes vacuous.
\item \textbf{Hyper‑parameter Sensitivity.} The bound depends linearly on $\eta_{\max}$ and inversely on $\eta_{\min}$. A larger amplitude $(\eta_{\max}-\eta_{\min})$ speeds up initial descent (larger bias term reduction) but inflates the steady‑state neighbourhood.
\item \textbf{Comparison to Baselines.} For standard (non‑cyclic) SGD with a constant step $\eta$, the classic bound reads $\frac{2(f(w_0)-f^{*})}{\eta T}+ \tfrac{L}{2}\eta\sigma^{2}$. Cyclic‑SGD replaces the constant $\eta$ by the *average* $\bar\eta=\frac{\eta_{\min}+\eta_{\max}}{2}$ in the bias term, while the variance term is governed by the worst‑case step $\eta_{\max}$. Hence cyclic schedules retain the same $O(1/T)$ decay of the bias term but may suffer a larger steady‑state error if $\eta_{\max}$ is much larger than the constant step used by vanilla SGD.
\end{itemize}

\section*{Discussion}
The theorem is derived directly from the smoothness inequality and the standard SGD analysis, the only novelty being the treatment of a *deterministic* but *periodic* step‑size sequence. Because the schedule is bounded, the proof proceeds without any averaging of the step sizes; the worst‑case step $\eta_{\max}$ appears naturally when bounding the variance contribution. When the cycle length $M$ grows with $T$ and the amplitudes shrink, the bound recovers the optimal $O(T^{-1/2})$ rate for non‑convex stochastic optimization.

\newpage
\section*{Preliminaries (Definitions and Lemmas)}
\begin{lemma}[One‑step Descent under Smoothness]\label{lem:descent}
For any $w\in\mathbb{R}^{d}$, any stochastic gradient $g$ and any step size $\eta>0$,
\[
f(w-\eta g)\le f(w)-\eta\langle\nabla f(w),g\rangle+\frac{L\eta^{2}}{2}\|g\|^{2}.
\tag{10}
\]
\end{lemma}
\begin{proof}
Apply smoothness inequality (4) with $u=w-\eta g$:
\[
f(w-\eta g)\le f(w)+\langle\nabla f(w),-\eta g\rangle+\frac{L}{2}\|\eta g\|^{2},
\]
which simplifies to (10). ∎
\end{proof}

\begin{lemma}[Variance Decomposition]\label{lem:variance}
For any random vector $g$ satisfying (5)–(6),
\[
\mathbb{E}\big[\|g\|^{2}\mid w\big]=\|\nabla f(w)\|^{2}+\mathbb{E}\big[\|g-\nabla f(w)\|^{2}\mid w\big]
\le \|\nabla f(w)\|^{2}+\sigma^{2}.
\tag{11}
\]
\end{lemma}
\begin{proof}
Write $g=\nabla f(w)+(g-\nabla f(w))$, expand the squared norm, and use unbiasedness (5) to kill the cross term. The inequality follows from (6). ∎
\end{proof}

\section*{Main Proof (Step‑by‑Step Derivation)}
We prove Theorem~\ref{thm:main}. Let $w_t$ be the iterate generated by (2) with step size $\eta_t$ defined by (1).

\noindent\textbf{[Step 1]} Apply Lemma~\ref{lem:descent} with $w=w_t$, $g=g_t$, $\eta=\eta_t$:
\[
f(w_{t+1})\le f(w_t)-\eta_t\langle\nabla f(w_t),g_t\rangle+\frac{L\eta_t^{2}}{2}\|g_t\|^{2}.
\tag{12}
\]

\noindent\textbf{[Step 2]} Take conditional expectation w.r.t. $g_t$ given $w_t$.
Using unbiasedness (5),
\[
\mathbb{E}\big[\langle\nabla f(w_t),g_t\rangle\mid w_t\big]
    =\|\nabla f(w_t)\|^{2}.
\tag{13}
\]

\noindent\textbf{[Step 3]} Apply Lemma~\ref{lem:variance} to bound the second moment term:
\[
\mathbb{E}\big[\|g_t\|^{2}\mid w_t\big]
    \le \|\nabla f(w_t)\|^{2}+\sigma^{2}.
\tag{14}
\]

\noindent\textbf{[Step 4]} Substitute (13)–(14) into (12) and then take total expectation $\mathbb{E}[\cdot]$:
\[
\mathbb{E}[f(w_{t+1})]
\le \mathbb{E}[f(w_t)]
      -\eta_t\,\mathbb{E}\!\big[\|\nabla f(w_t)\|^{2}\big]
      +\frac{L\eta_t^{2}}{2}\,\mathbb{E}\!\big[\|\nabla f(w_t)\|^{2}+\sigma^{2}\big].
\tag{15}
\]

\noindent\textbf{[Step 5]} Rearrange (15) to isolate the gradient term:
\[
\Bigl(\eta_t-\frac{L\eta_t^{2}}{2}\Bigr)\mathbb{E}\!\big[\|\nabla f(w_t)\|^{2}\big]
\le \mathbb{E}[f(w_t)]-\mathbb{E}[f(w_{t+1})]+\frac{L\eta_t^{2}}{2}\sigma^{2}.
\tag{16}
\]

\noindent\textbf{[Step 6]} Since $0<\eta_t\le\eta_{\max}$, we have
\[
\eta_t-\frac{L\eta_t^{2}}{2}
\;\ge\;
\eta_{\min}-\frac{L\eta_{\max}^{2}}{2}
\;\ge\;
\frac{\eta_{\min}}{2},
\tag{17}
\]
provided the steps satisfy $L\eta_{\max}<1$ (a standard stability condition).  Hence
\[
\frac{\eta_{\min}}{2}\,\mathbb{E}\!\big[\|\nabla f(w_t)\|^{2}\big]
\le
\mathbb{E}[f(w_t)]-\mathbb{E}[f(w_{t+1})]+\frac{L\eta_{\max}^{2}}{2}\sigma^{2}.
\tag{18}
\]

\noindent\textbf{[Step 7]} Sum inequality (18) over $t=0,\dots,T-1$:
\[
\frac{\eta_{\min}}{2}\sum_{t=0}^{T-1}\mathbb{E}\!\big[\|\nabla f(w_t)\|^{2}\big]
\le
f(w_{0})- \mathbb{E}[f(w_{T})] + \frac{L\eta_{\max}^{2}}{2}\,T\sigma^{2}.
\tag{19}
\]

\noindent\textbf{[Step 8]} Drop the non‑negative term $-\mathbb{E}[f(w_T)]\le -f^{*}$ and divide by $T$:
\[
\frac{1}{T}\sum_{t=0}^{T-1}\mathbb{E}\!\big[\|\nabla f(w_t)\|^{2}\big]
\le
\frac{2\big(f(w_{0})-f^{*}\big)}{\eta_{\min}T}
   + L\,\eta_{\max}\,\sigma^{2}.
\tag{20}
\]
Equation (20) is exactly the bound (9) claimed in Theorem~\ref{thm:main}. ∎

\section*{Rate Derivation (With Constants)}
From (20) we identify two contributions:
\begin{itemize}
\item \textbf{Bias decay term} $ \displaystyle \frac{2\big(f(w_{0})-f^{*}\big)}{\eta_{\min}T}$, which decays as $O(1/T)$.
\item \textbf{Variance neighbourhood term} $L\eta_{\max}\sigma^{2}$, independent of $T$.
\end{itemize}
If we let the amplitude shrink with $T$ by setting $\eta_{\max}=c/\sqrt{T}$ (and similarly $\eta_{\min}=c'/\sqrt{T}$) then the variance term becomes $L c\sigma^{2}/\sqrt{T}$ and the whole right‑hand side scales as $O(T^{-1/2})$, which matches the optimal rate for non‑convex stochastic optimization.

\section*{Special Cases}
\begin{enumerate}
\item \textbf{Constant Cyclic Schedule.} When $\eta_{\min}$ and $\eta_{\max}$ are fixed, (20) shows convergence to a neighbourhood of size $L\eta_{\max}\sigma^{2}$. If the stochastic noise vanishes ($\sigma=0$) the neighbourhood collapses and the algorithm attains exact stationarity at rate $O(1/T)$.
\item \textbf{Deterministic Gradient (Full‑batch).} Setting $\sigma=0$ in (20) recovers the deterministic bound
\[
\frac{1}{T}\sum_{t=0}^{T-1}\|\nabla f(w_t)\|^{2}
\le \frac{2\big(f(w_{0})-f^{*}\big)}{\eta_{\min}T},
\]
i.e. the average gradient norm decays as $O(1/T)$ irrespective of the cyclic modulation.
\item \textbf{Large Cycle Length $M$.} As $M\to\infty$ the schedule (1) becomes effectively constant $\eta_t\approx\eta_{\min}$, and the bound reduces to the classic SGD bound with step $\eta_{\min}$.
\end{enumerate}

\section*{Conclusion}
We have presented a full, step‑by‑step convergence analysis for Stochastic Gradient Descent equipped with a triangular cyclic learning‑rate schedule. The proof follows directly from smoothness and unbiasedness assumptions, with explicit handling of the periodic step‑size sequence. The resulting bound (9) captures both the $O(1/T)$ decay of the optimization bias and the steady‑state error dictated by the maximal learning rate and gradient variance, thereby providing a rigorous theoretical foundation for the empirical success of cyclic learning‑rate strategies.

\end{document}